\section{Cinemática e modelagem de manipuladores espaciais}

A modelagem cinemática de manipuladores espaciais constitui uma etapa central para descrever o movimento relativo entre os elos, determinar a posição e a orientação do efetuador final e estabelecer as relações entre o espaço das juntas e o espaço cartesiano \cite{fonseca2019kinematics}. No contexto espacial, essa modelagem adquire complexidade adicional, pois o manipulador está acoplado a uma plataforma que não pode ser tratada, em geral, como uma base fixa convencional. Em vez disso, a formulação precisa considerar a interação entre o manipulador e a espaçonave, bem como o uso consistente de referenciais inercial, orbital e fixo ao corpo \cite{fonseca2019kinematics,fonseca2016dynamics}.

Nesse cenário, abordagens clássicas de cinemática direta e inversa continuam desempenhando papel fundamental. A convenção de Denavit--Hartenberg permanece amplamente empregada para a definição sistemática dos referenciais e para a obtenção das transformações homogêneas entre os elos, permitindo calcular a pose do efetuador final a partir das variáveis articulares \cite{fonseca2019kinematics}. Quando a cinemática inversa não admite solução analítica em forma fechada, o problema pode ser formulado como uma otimização, em que se busca a configuração articular que minimiza o erro entre a pose desejada e a pose obtida pelo manipulador \cite{santos2022uncertainty,santos2023physics}.

Entretanto, em manipuladores do tipo espaçonave, a formulação cinemática não se limita à descrição geométrica da cadeia serial. Como a movimentação das juntas pode induzir deslocamentos e rotações da plataforma, torna-se necessário compatibilizar a cinemática do manipulador com os movimentos da espaçonave e com os referenciais relevantes ao problema orbital \cite{fonseca2019kinematics,fonseca2016dynamics}. Essa característica é especialmente importante em operações de proximidade, captura e serviço em órbita, nas quais erros de posicionamento e orientação podem comprometer o desempenho da missão \cite{arantes2010far,fonseca2016dynamics}.

Além das formulações analíticas tradicionais, estudos recentes têm explorado estratégias baseadas em aprendizado de máquina para apoiar a resolução da cinemática inversa e o planejamento de movimento. Redes neurais e métodos como máquinas de vetores de suporte vêm sendo empregados para estimar configurações articulares, quantificar incertezas e melhorar a convergência de procedimentos de otimização, inclusive em situações em que os efeitos dinâmicos da plataforma e do manipulador são incorporados ao problema \cite{santos2022uncertainty,santos2023physics,santos2022machine}. Esses resultados indicam uma tendência de combinar modelos cinemáticos clássicos com técnicas computacionais orientadas por dados, visando ampliar a autonomia e a robustez operacional de manipuladores espaciais.

De modo geral, a literatura apresenta contribuições consistentes para a cinemática de manipuladores espaciais, tanto sob a perspectiva geométrica tradicional quanto sob abordagens computacionais mais recentes. Ainda assim, observa-se que parte dos trabalhos concentra-se na determinação da pose e na resolução da cinemática inversa, enquanto outra parte busca incorporar, de forma mais explícita, os efeitos do acoplamento entre manipulador e plataforma. Essa distinção é relevante para a presente dissertação, pois a modelagem cinemática constitui a base para a formulação posterior da dinâmica acoplada do sistema durante operações orbitais de proximidade \cite{fonseca2019kinematics,fonseca2016dynamics,santos2022uncertainty}.